% New English translation from checked cote 25, site 99 (PDF page 100).
% AI-translated by GPT-6 Astra (OpenAI, ChatGPT Pro mode).
\scanpage{25/99}
% NOTE: the reviewed source records two long diagonal pen strokes; their function does not establish cancellation of the entire argument.
\mspage{2}

obtained from $f:P\hand{\longrightarrow} S$, the second determined by the element cl $(\underline{O}_P(1))$ of
Pic(P). We infer a canonical homomorphism
(x) \qquad Pic(S)x\var{\xout{\unsure{ZZx(}}}{}{typewriter strike-out}{ega25_25_099-v0001} $\underline{Z} \hand{\longrightarrow}$ Pic(P),
defined moreover without any restrictive hypothesis on S or on $\underline{E}$. This said:
\uline{Corollary} 1.4.\var{\margin{\ensuremath{\swarrow}}}{\placedmargin{\ensuremath{\swarrow}}}{marginal note placed}{ega25_25_099-v0002} Under the conditions of 1.\unsure{1|4}., \var{\inserted{if $S\neq\emptyset$}}{if $S\neq\emptyset$}{insertion}{ega25_25_099-v0003} \var{\xout{suppose S connected non}}{}{typewriter strike-out}{ega25_25_099-v0004}
\var{\xout{empty}}{}{typewriter strike-out}{ega25_25_099-v0005} then the preceding homomorphism is \var{\xout{an isomorphism}}{}{typewriter strike-out}{ega25_25_099-v0006} \var{\inserted[typed]{injective}}{injective}{insertion}{ega25_25_099-v0007}, and even
bijective when S is connected\hand{.}\var{\deleted{nonempty.}}{}{deletion}{ega25_25_099-v0008}

\var{\margin{[}}{\placedmargin{[}}{marginal note placed}{ega25_25_099-v0009} When we \var{\xout{suppose only S connected, and}}{}{typewriter strike-out}{ega25_25_099-v0010} abandon the hypothesis
on the rank $\hand{\geqslant} 2$, it follows from 1.2. and 1.\unsure{4|3}. that
\var{the preceding homomorphism \inserted[typed]{if S is connected, but}
is still surjective}{the preceding homomorphism is still surjective if S is connected, but}{insertion}{ega25_25_099-v0011} \var{\xout{(but}}{}{typewriter strike-out}{ega25_25_099-v0012} not necessarily \unsure{injective|bijective}, the kernel being
isomorphic \var{\xout{to Z resp. to Pic(S)xZ when E is of rank 1 resp. of rank 0)}}{}{typewriter strike-out}{ega25_25_099-v0013}
to $\underline{Z}$ resp. to Pic(S)x$\underline{Z}$ when $\underline{E}$ is of rank 1\hand{,} resp. of rank 0.

Let us prove 1.1. beginning with uniqueness. First, if S is
the spectrum of a field, note that $\underline{O}(n)$ is isomorphic to $\underline{O}(m)$ i.e.
$\underline{O}(m-n)$ is isomorphic to $\underline{O}_P$ only if n=m i.e. n-m=0. This follows
from the fact that $\underline{O}(1)$ is ample and dim P $\hand{\geqslant}$1 (by the hypothesis rank E $\hand{\geqslant}$2):
we can indeed suppose d = m-n $\hand{\geqslant}$0; if we had d $\hand{>}$0, then $\underline{O}(d)$
would be ample, and could be isomorphic to $\underline{O}_P$ only if P is quasi-affine,
hence finite (since it is proper over k). This already proves uniqueness
of the family $(S_n)_{n\hand{\in}\underline{Z}}$ considered in 1.1. For uniqueness of $\underline{M}$ up to unique
isom., we reduce to the case where $S=S_n$; I say that in this case we have
an isomorphism (uniquely determined in terms of the isom $\underline{L} \hand{\xrightarrow{\sim}} f$\var{\over{$_*$}{$^*$}}{$^*$}{writing over}{ega25_25_099-v0014}$(\underline{M})(n)$):
\[ \text{(xx)} \qquad M \hand{\xrightarrow{\sim}} f_*(\underline{\unsure{L|E|M}}(-n)) \quad . \]
Indeed, the isomorphism $\underline{L} \hand{\xrightarrow{\sim}} f^*(\underline{M})(n)$ defines an isomorphism
$\underline{L}(-n) \hand{\xrightarrow{\sim}} f^*(\underline{M})$ and consequently an isomorphism of the second member of (xx)
with $f_*(f^*(\underline{M}))$, which itself is isomorphic \unsure{(|to}\,$\underline{M}$ being loc free) to
$\underline{M}\otimes f_*(\underline{O}_P)$, but (ref) $f_*(\underline{O}_P) \hand{\xleftarrow{\sim}} \underline{O}_{\unsure{S|S_S}}$, whence the \unsure{isomorphism|isomorphism} (xx).\var{\margin{\ensuremath{\backslash}}}{\placedmargin{\ensuremath{\backslash}}}{marginal note placed}{ega25_25_099-v0015}
\tnote{25-099-1}{The two spellings in the final live uncertainty are isomorphisme and isemorphisme; they give the same English word. The source's plain M in (xx), its doubtful sheaf letter, the unsupplied reference ref, and the underlined O subscript alternatives are retained.}

\var{\margin{93}}{\placedmargin{93}}{marginal note placed}{ega25_25_099-v0016}
